\section*{Введение}
\addcontentsline{toc}{section}{Введение}

\textbf{Актуальность темы.} В современном мире информационные технологии играют ключевую роль в развитии общества, а навыки программирования становятся необходимыми не только для специалистов в области IT, но и для широкого круга пользователей. Особую значимость приобретает подготовка к олимпиадам по программированию, которые развивают алгоритмическое мышление, способность решать нестандартные задачи и повышают уровень профессиональной подготовки студентов. Однако существующие образовательные ресурсы зачастую либо перегружены теорией, либо не предоставляют удобных инструментов для практической отработки навыков и анализа прогресса. В связи с этим возникает необходимость создания современного веб-приложения, которое объединяет обучение, практику и контроль знаний в единой среде.

\textbf{Степень разработанности проблемы.} В настоящее время существует множество платформ для изучения программирования. Тем не менее, многие из них ориентированы либо на начальный уровень подготовки, либо на узкоспециализированные задачи. Недостатком большинства решений является отсутствие персонализации обучения, ограниченные возможности анализа ответов пользователя и слабая интеграция инструментов подготовки к олимпиадам. Кроме того, не все платформы обеспечивают удобное взаимодействие между пользователем и системой, включая гибкую проверку решений и рекомендации по улучшению.

\textbf{Объект исследования.} Объектом исследования является процесс обучения программированию и подготовки к олимпиадам в условиях использования современных веб-технологий.


\textbf{Предмет исследования.} Предметом исследования выступают методы и средства разработки веб-приложения, обеспечивающего эффективное изучение программирования, включая хранение задач, автоматическую проверку решений, анализ ответов пользователя и предоставление обратной связи.

\textbf{Целью данной работы} является разработка веб-приложения для изучения программирования и подготовки к олимпиадам, которое обеспечивает удобный интерфейс, систему хранения и выдачи задач, интеллектуальную проверку ответов пользователей и анализ их уровня знаний с последующей выдачей рекомендаций.